\section{Анализ сходимости}
\label{ch:theory}

Особенность Алгоритма SGEM в том, что суррогат $u^{(t)}$, используемый на
E- и M-шагах, сам обновляется между итерациями, поэтому классический
аргумент монотонности EM~\cite{dempster1977em, wu1983convergence}
напрямую неприменим. Покажем, что при достаточно быстром убывании ошибки
суррогата SGEM сходится в том же смысле, что и обобщённый EM.

Пусть $u_{\mathrm{true}}$ --- неизвестная истинная ошибка эксперта
(пообъектная кросс-энтропия архитектуры $\boldsymbol{\alpha}$, обученной
на подмножестве $\mathbf{r}$). Определим пошаговую ошибку суррогата в
шкале потерь
\begin{equation}
\label{eq:eps_t}
\varepsilon_t \;=\; \sup_{\boldsymbol{\alpha} \in \mathcal{A}, \;
\mathbf{r} \in [0, 1]^M}
\bigl|u^{(t)}(\boldsymbol{\alpha}, \mathbf{r})
- u_{\mathrm{true}}(\boldsymbol{\alpha}, \mathbf{r})\bigr|.
\end{equation}
Параметр $\theta = (\boldsymbol{\alpha}_{1:K}, \mathbf{R})$ лежит в
компактном множестве $\Theta = \mathcal{A}^K \times (\Delta_K)^M$;
обозначим $N = \sum_m |\mathcal{C}_m|$. Истинное (взвешенное по размерам
кластеров) лог-правдоподобие есть
\begin{equation}
\label{eq:Ltrue}
L_{\mathrm{true}}(\theta) \;=\; \sum_{m=1}^M |\mathcal{C}_m|\,\log
\sum_{k=1}^K r_{mk}\, e^{-u_{\mathrm{true}}(\boldsymbol{\alpha}_k, \mathbf{r}_k)},
\end{equation}
а истинная $Q$-функция ---
\begin{equation}
\label{eq:Qtrue}
Q_{\mathrm{true}}(\theta; \theta') \;=\; \sum_{m=1}^M |\mathcal{C}_m|
\sum_{k=1}^K w_{mk}(\theta')
\bigl[\log r_{mk} - u_{\mathrm{true}}(\boldsymbol{\alpha}_k, \mathbf{r}_k)\bigr],
\end{equation}
где $w_{mk}(\theta')$ --- ответственности, вычисленные при
$u_{\mathrm{true}}$. Введём вариационную свободную энергию,
параметризованную функцией потерь $u$:
\[
\mathcal{F}_u(q, \theta) \;=\; \sum_{m=1}^{M} |\mathcal{C}_m|
\sum_{k=1}^{K} q_{mk}\, \log \frac{r_{mk}\, e^{-u(\boldsymbol{\alpha}_k, \mathbf{r}_k)}}{q_{mk}},
\qquad L_u(\theta) = \max_q \mathcal{F}_u(q, \theta).
\]

\begin{assumption}
\label{ass:bounded}
Существует $U_{\max} < \infty$ такое, что
$0 \leq u_{\mathrm{true}}, u^{(t)} \leq U_{\max}$ для всех
$\boldsymbol{\alpha}, \mathbf{r}, t$.
\end{assumption}

\begin{assumption}
\label{ass:summable}
Ошибки суррогата суммируемы:
$\sum_{t=0}^{\infty} \varepsilon_t < \infty$.
\end{assumption}

\begin{assumption}
\label{ass:gem}
M-шаг обобщён, т.е. для каждого $t$ возвращается $\theta^{(t+1)}$ с
$\mathcal{F}_{u^{(t)}}(q^{(t)}, \theta^{(t+1)})
\geq \mathcal{F}_{u^{(t)}}(q^{(t)}, \theta^{(t)})$.
\end{assumption}

\begin{assumption}
\label{ass:cont}
Для каждого $\boldsymbol{\alpha}$ отображение
$\mathbf{r}\mapsto u_{\mathrm{true}}(\boldsymbol{\alpha}, \mathbf{r})$
непрерывно на $[0,1]^M$.
\end{assumption}

\begin{theorem}[Сходимость SGEM с адаптивным суррогатом]
\label{thm:sgem}
При выполнении Предположений~\ref{ass:bounded}--\ref{ass:cont}
последовательность $\{\theta^{(t)}\}$, порождаемая Алгоритмом SGEM,
удовлетворяет:
\begin{enumerate}
\item[\textup{(a)}] последовательность
$\{L_{\mathrm{true}}(\theta^{(t)})\}$ сходится;
\item[\textup{(b)}] всякая предельная точка $\theta^{\star}$ является
стационарной точкой $L_{\mathrm{true}}$ на $\Theta$ в смысле
$Q_{\mathrm{true}}(\theta;\theta^{\star}) \leq
Q_{\mathrm{true}}(\theta^{\star};\theta^{\star})$ для всех
$\theta\in\Theta$.
\end{enumerate}
\end{theorem}

\begin{proof}[Схема доказательства]
Ключевые шаги следующие. Из вариационного представления
$L_u(\theta) = \max_q \mathcal{F}_u(q, \theta)$ с максимумом в точке
$q^{\star}_{mk} \propto r_{mk} e^{-u(\boldsymbol{\alpha}_k, \mathbf{r}_k)}$
следует равномерная оценка отклонения:
$|\mathcal{F}_{u_{\mathrm{true}}}(q,\theta) - \mathcal{F}_{u^{(t)}}(q,\theta)|
\leq N\varepsilon_t$ и, как следствие,
$|L_{\mathrm{true}}(\theta) - L^{(t)}(\theta)| \leq N\varepsilon_t$.
Объединяя это с Предположением~\ref{ass:gem} (подъём свободной энергии на
M-шаге) и равенством на E-шаге, получаем
\emph{квазимонотонность}
\begin{equation}
\label{eq:quasi-mono}
L_{\mathrm{true}}(\theta^{(t+1)}) \;\geq\;
L_{\mathrm{true}}(\theta^{(t)}) \;-\; 2N\,\varepsilon_t.
\end{equation}
Положив $V_t = L_{\max} - L_{\mathrm{true}}(\theta^{(t)}) \geq 0$
(величина $L_{\max} = \sup_\theta L_{\mathrm{true}}$ конечна ввиду
компактности $\Theta$), из~\eqref{eq:quasi-mono} имеем
$V_{t+1} \leq V_t + 2N\varepsilon_t$ при
$\sum_t 2N\varepsilon_t < \infty$. По теореме Роббинса--Зигмунда
\cite{robbins1971convergence} последовательность $\{V_t\}$ сходится, что
доказывает~(a). Для~(b) используется устойчивость softmax-ответственностей
к возмущению суррогата: при $\varepsilon_t \to 0$ ответственности и
$Q$-функция, вычисленные при $u^{(t)}$, сходятся к истинным; если бы
предельная точка не была стационарной, нашёлся бы $\theta'$ с
$Q_{\mathrm{true}}(\theta';\theta^{\star}) -
Q_{\mathrm{true}}(\theta^{\star};\theta^{\star}) = \delta > 0$, что по
переходу через $\tilde Q^{(t)}$ и подъёму на M-шаге дало бы прирост
$\Delta_t \geq \delta/8$ для бесконечного числа $t$, противоречащий
суммируемости из~(a). Следовательно, $\theta^{\star}$ стационарна.
\end{proof}

Побочным следствием доказательства пункта~(a) является явная оценка
скорости
\[
L_{\mathrm{true}}(\theta^{\star}) - L_{\mathrm{true}}(\theta^{(T)})
\;\leq\; 2N \sum_{t = T}^{\infty} \varepsilon_t.
\]
Если $\varepsilon_t = O(t^{-p})$ при $p > 1$, то
$L_{\mathrm{true}}(\theta^{(T)}) \to L_{\mathrm{true}}(\theta^{\star})$
со скоростью $O(T^{1 - p})$, что согласуется со скоростями,
наблюдаемыми в анализе стохастического EM~\cite{cappe2009online}.

\begin{remark}
Предположение суммируемости~\ref{ass:summable} --- суррогатный аналог
условия убывающего шага в стохастическом EM; ослабление его до
$\varepsilon_t \to 0$ сохраняет пункт~(b), но не пункт~(a). На практике
суммируемость достигается, когда шаг активного обучения обеспечивает
достаточно полное покрытие пространства
$\mathcal{A}\times[0,1]^M$.
\end{remark}

\textbf{Новизна и связь с классическими результатами о сходимости.}
Классическая теория сходимости EM и его обобщённого варианта
\cite{dempster1977em, wu1983convergence} опирается на монотонный подъём
единого \emph{фиксированного} правдоподобия: каждый M-шаг не уменьшает
одну и ту же целевую функцию, и сходимость следует из монотонной
ограниченной последовательности. В нашей постановке эта предпосылка
нарушена. Поскольку качество эксперта берётся с адаптивно
переобучаемого суррогата $u^{(t)}$, а не из точного обучения,
оптимизируемая на итерации $t$ цель сама является «движущейся мишенью»,
а истинное правдоподобие $L_{\mathrm{true}}$, которое нас в итоге
интересует, алгоритмом \emph{никогда} не вычисляется. Поэтому точная
монотонность недоступна, и стандартный аргумент неприменим.

Теорема~\ref{thm:sgem} показывает, что сходимость тем не менее
сохраняется, --- в этом состоит основной теоретический вклад. Механизм
--- \emph{квази}-монотонность
$L_{\mathrm{true}}(\theta^{(t+1)}) \geq L_{\mathrm{true}}(\theta^{(t)}) - 2N\varepsilon_t$:
истинное правдоподобие может проседать, но лишь на величину,
пропорциональную пошаговой ошибке суррогата $\varepsilon_t$.
Суммируемость этих ошибок превращает
$\{L_{\mathrm{true}}(\theta^{(t)})\}$ в почти-супермартингал, и аргумент
Роббинса--Зигмунда~\cite{robbins1971convergence} даёт и сходимость, и
стационарность всякой предельной точки \emph{истинной} цели --- хотя
каждый E- и M-шаг видит только суррогат.

Это отличается от анализа стохастического и онлайн-EM
\cite{cappe2009online}, где возмущение точного шага EM ---
\emph{нулевое в среднем стохастическое} (от Монте-Карло E-шага),
гасимое убывающим шагом. В нашем случае возмущение --- наоборот,
\emph{детерминированная ограниченная ошибка аппроксимации} обученного
суррогата, а роль убывающего шага играет условие
суммируемости~\ref{ass:summable}, малость которого обеспечивается не
усреднением шума, а покрытием пространства поиска при активном обучении.
Насколько нам известно, это первая гарантия сходимости для EM-подобной
процедуры архитектурного поиска, чья пошаговая цель задаётся
\emph{адаптивно обучаемой} суррогатной моделью.
